\section{Présentation de la méthode}

Émile Zola, l'un des écrivains majeurs du XIXe siècle, est connu pour ses romans naturalistes, dans lesquels il explore les dimensions sociales, économiques et psychologiques de la vie humaine. Son œuvre romanesque constitue ainsi un corpus particulièrement riche pour l'analyse des thèmes récurrents, des milieux sociaux et des motifs narratifs.

Dans cette étude, nous examinons les thèmes présents dans les romans de Zola afin de déterminer dans quelle mesure certains motifs apparaissent de manière transversale dans son œuvre. Pour cela, nous mobilisons des méthodes d'analyse textuelle non supervisées, capables d'identifier des regroupements lexicaux à partir des textes eux-mêmes.

La question centrale de cette étude est donc la suivante : dans quelle mesure des méthodes non supervisées comme la NMF et K-means permettent-elles de faire apparaître des structures thématiques récurrentes dans l'œuvre romanesque de Zola ?

L'objectif est de vérifier si les regroupements produits automatiquement correspondent à des motifs interprétables dans le cadre d'une analyse littéraire. Il ne s'agit donc pas de remplacer la lecture critique des textes, mais d'utiliser les outils du traitement automatique du langage comme un moyen d'explorer un corpus littéraire volumineux.

Ce travail constitue le premier volet d'une série d'analyses portant sur des corpus de textes littéraires issus d'auteurs naturalistes et non naturalistes. Il sert à la fois de cas d'étude méthodologique et de base pour de futures comparaisons entre différents ensembles d'œuvres.

Deux méthodes d'analyse ont été utilisées pour identifier les thèmes dans les romans de Zola :

\begin{itemize}
    \item la NMF (Non-negative Matrix Factorization)\footnote{N. Gillis, \textit{Nonnegative Matrix Factorization}, Society for Industrial and Applied Mathematics, coll. « Data Science », 2020.} : une technique de réduction de dimensionnalité qui permet d'extraire des topics à partir d'un corpus de textes ;
    \item K-means : une méthode de clustering qui permet de regrouper les segments textuels en fonction de leur proximité lexicale.
\end{itemize}

Ces méthodes seront expliquées dans les sections suivantes, avant la présentation des résultats de l'analyse thématique des romans de Zola.

Dans cette étude, le terme \textit{topic} désigne un regroupement lexical produit par le modèle NMF, tandis que le terme \textit{cluster} désigne un groupe de segments obtenu par K-means. Le terme \textit{thème} est réservé à l'interprétation qualitative de ces regroupements. Cette précision terminologique est importante pour éviter toute confusion entre les différentes étapes de l'analyse.